\documentclass[11pt,a4paper]{article}
\usepackage[T1]{fontenc}
\usepackage[utf8]{inputenc}
\usepackage{lmodern}
\usepackage{amsmath,amssymb,amsthm,mathtools}
\usepackage{esint,graphicx}
\usepackage[a4paper,margin=22mm,headheight=14pt]{geometry}
\usepackage{microtype}
\usepackage{booktabs,longtable,array,enumitem}
\newcolumntype{L}[1]{>{\raggedright\arraybackslash}p{#1}}
\usepackage[dvipsnames]{xcolor}
\usepackage{fancyhdr}
\usepackage{hyperref}

\newtheorem{theorem}{Theorem}[section]
\newtheorem{lemma}[theorem]{Lemma}
\newtheorem{proposition}[theorem]{Proposition}
\newtheorem{corollary}[theorem]{Corollary}
\theoremstyle{remark}
\newtheorem{remark}[theorem]{Remark}
\newtheorem{definition}[theorem]{Definition}
\newtheorem{claim}[theorem]{Claim}
\setlist{nosep,leftmargin=1.5em}
\setlength{\emergencystretch}{2em}
\pagestyle{fancy}
\fancyhf{}
\fancyhead[L]{\small AIM 721}
\fancyhead[R]{\small Revised 7 November 2026}
\fancyfoot[C]{\thepage}
\renewcommand{\headrulewidth}{0.3pt}
\hypersetup{colorlinks=true,linkcolor=MidnightBlue,urlcolor=MidnightBlue,citecolor=MidnightBlue,pdftitle={AIM 721},pdfauthor={Matthew J. Colbrook; results found by AI}}
\begin{document}
\title{Real regularity of complex Hessian obstacle envelopes}
\author{Matthew J. Colbrook\thanks{The solutions and partial results in this document were found by AI. The supplied research and the separate repository review are AI work; human peer review and formal verification are not claimed.}\\[4pt]
\small Department of Applied Mathematics and Theoretical Physics\\
\small University of Cambridge, Cambridge CB3 0WA, United Kingdom\\
\small\href{mailto:m.colbrook@damtp.cam.ac.uk}{m.colbrook@damtp.cam.ac.uk}}
\date{7 November 2026}
\maketitle
\begin{abstract}
The complex $m$-Hessian obstacle envelope is locally $C^{1,1}$ in all real coordinates for every $1\le m\le n$, under the stated strict defining-function assumption. The penalized Bellman estimates are uniform in both approximation parameters. The conclusion concerns exactly the catalogue target; no historical priority claim is made.
\end{abstract}
\noindent\textbf{Catalogue scope.} AIM 721; complete resolution. \href{https://github.com/MColbrook/AIM/blob/1e564633494c109edfd3ec46da1366adcff925a5/problems/721-real-c1-1-regularity-of-complex-hessian-obstacle-envelopes.md}{Canonical target at revision \texttt{1e56463}}.
The domain-envelope question is asked in \cite[p.~998]{BenaliZeriahi}. The accompanying erratum \cite{BenaliErratum} is retained in the references; its incomplete original boundary-data argument is not used here.

\section{Statement and argument}
\label{sec:721}

The exact obstacle-envelope question has an affirmative answer under the
stated strict defining-function hypothesis.  The argument uses smooth Bellman
approximants with a quadratic stopping penalty.  The essential new estimate
compensates the discount when the diffusion clock is perturbed; this keeps the
large discount out of the second-variation bound.  We give the approximation,
control comparison, and limiting argument explicitly.

\begin{theorem}
\label{thm:721}
Let $\Omega\Subset\mathbb C^n$ have smooth boundary and a smooth defining
function $\rho$ on a neighborhood of $\overline\Omega$ such that
$\rho<0$ in $\Omega$, $d\rho\ne0$ on $\partial\Omega$, and $\rho$ is
strictly $m$-subharmonic on that neighborhood.  If
$h\in C^2(\overline\Omega;\mathbb R)$, then
\begin{equation}
 P_{m,\Omega}h
 :=\bigl(\sup\{w\in\operatorname{SH}_m(\Omega):w\le h\}\bigr)^*
 \in C^{1,1}_{\mathrm{loc}}(\Omega)\cap C^{0,1}(\overline\Omega).
 \label{eq:721-1}
\end{equation}
Here $C^{1,1}$ concerns the full real Hessian.  The conclusion holds for all
$1\le m\le n$, including the intermediate orders in AIM 721.
\end{theorem}

\subsection{Covariances and smooth approximants}

Let $\Gamma_m\subset\operatorname{Herm}(n)$ be the open G\aa rding cone,
and put
\begin{equation}
 \mathcal A=\{A\in\overline\Gamma_m^{\,*}:\operatorname{tr}A=1\},
 \qquad L_Au=\operatorname{tr}(Au_{z\bar z}).
 \label{eq:721-2}
\end{equation}
The cone contains the positive definite cone, so its dual consists of positive
semidefinite matrices.  Consequently $\mathcal A$ is compact and convex.
It is invariant under unitary conjugation and contains $I/n$, since
$\Gamma_m\subset\{H:\operatorname{tr}H>0\}$.  The constant-coefficient
dual-cone characterization says that $u$ is $m$-subharmonic if and only if
$L_Au\ge0$ for every $A\in\mathcal A$, in the distributional, equivalently
viscosity, sense.  For smooth functions this is convex duality; the general
statement follows by convolution and the usual viscosity characterization of
constant-coefficient linear subharmonicity.

Strictness and compactness give $L_A\rho\ge\delta_0>0$ uniformly on
$\overline\Omega$ and $A\in\mathcal A$.  Scale $-\rho$ to obtain $\psi$
with
\begin{equation}
 \Omega=\{\psi>0\},\qquad \psi=0\text{ on }\partial\Omega,
 \qquad |\nabla\psi|\ge1\text{ on }\partial\Omega,
 \qquad L_A\psi\le-1.
 \label{eq:721-3}
\end{equation}
Enlarge this fixed scaling, if necessary, so that
$-L_A\psi\ge4\sup_A\|\sigma_A\|^2$, where in real coordinates
$L_A=(1/2)\sigma_A\sigma_A^{\mathsf T}:D^2$.  This is the normalization
used immediately before \cite[Lemma 7.1]{ZhouBellman2015}.
All covariance traces are bounded uniformly.

Set $g=-h$ and, for $0<\nu<1$, define
\begin{equation}
 \mathcal A_\nu=\{(1-\nu)A+\nu I/n:A\in\mathcal A\},\qquad
 F_\nu(D^2v)=\sup_{A\in\mathcal A_\nu}L_Av.
 \label{eq:721-4}
\end{equation}
The sets $\mathcal A_\nu$ are unitary invariant, contain $I/n$, and satisfy
\eqref{eq:721-3}.  Each fixed $\nu$ gives a uniformly elliptic operator.  The estimates
below use an upper bound for the covariances but no positive lower bound.
Initially take $g$ smooth on a neighborhood of $\overline\Omega$.

Choose
\[
 C_0\ge\sup_{\nu,A\in\mathcal A_\nu}\|L_Ag\|_\infty,
 \qquad C_*=C_0+1,\qquad b_K=\sqrt{2C_*K},\qquad K\ge1,
\]
and write
\[
 H_b(s)=\max_{0\le c\le b}\left(cs-\frac{c^2}{2K}\right).
\]
For fixed $K,\nu$ solve
\begin{equation}
 F_\nu(D^2v)+H_{2b_K}(g-v)=0\quad\text{in }\Omega,
 \qquad v=g\quad\text{on }\partial\Omega.
 \label{eq:721-5}
\end{equation}
This is a uniformly elliptic Bellman equation with compact controls
$(A,c)\in\mathcal A_\nu\times[0,2b_K]$ and running payoff
$cg-c^2/(2K)$.  Standard nondegenerate Bellman Dirichlet theory gives
$v=v_{K,\nu}\in C^{2,\alpha}(\overline\Omega)$ for some $\alpha>0$.
Precisely this fixed-parameter regularization is used in
\cite[\S5, Remark 5.1]{ZhouBellman2015}, citing
\cite[\S6.2, Theorem 7]{Krylov1987}.  Only the existence of a $C^2$ solution
for fixed parameters is used here, not a uniform estimate from that theorem.

\paragraph{The two discount caps.}
At a positive maximum of $w=g-v$, one has $D^2v\ge D^2g$, whence
\[
 H_{2b_K}(w)=-F_\nu(D^2v)\le-F_\nu(D^2g)\le C_0.
\]
If $w\ge b_K/K$, the admissible choice $c=b_K$ would give
$H_{2b_K}(w)\ge C_*>C_0$.  Thus $w<b_K/K$ everywhere, the maximizing
discount is $c=K(g-v)_+$, and
\begin{equation}
 F_\nu(D^2v)+\frac K2(g-v)_+^2=0,
 \qquad (g-v)_+\le\sqrt{2C_0/K}.
 \label{eq:721-6}
\end{equation}
The same $v$ therefore solves the Bellman problem with the smaller cap
$[0,b_K]$, and verification identifies it with the value function for either
cap.  We shall use base controls $c_t\le b_K$, so that
\begin{equation}
 q_t:=\frac{c_t^2}{2K}\le C_*.
 \label{eq:721-7}
\end{equation}
A perturbed control will use the larger cap.  This bound holds for every base
control, without any assumption of exact optimality.

\paragraph{Uniform value and gradient bounds.}
For a fixed sufficiently large $C$, $g+C\psi$ and $g-C\psi$ are respectively
a supersolution and a subsolution of \eqref{eq:721-5}.  Indeed the first has
$F_\nu(D^2(g+C\psi))\le C_0-C$ and zero penalty, whereas the second has
$F_\nu(D^2(g-C\psi))\ge-C_0+C$ and nonnegative penalty.  Thus
\begin{equation}
 |v-g|\le C\psi.
 \label{eq:721-8}
\end{equation}
For a small real translation $a$, put
$v_a(x)=v(x+a)-L|a|$ on $\Omega\cap(\Omega-a)$.  If
$L\ge\operatorname{Lip}(g)$, then
$g(x)-v_a(x)\ge g(x+a)-v(x+a)$, so $v_a$ is a subsolution of \eqref{eq:721-5}.
On the boundary of the intersection, one of $x,x+a$ lies on
$\partial\Omega$.  Formula \eqref{eq:721-8} gives $v_a\le v$ there if
$L\ge\operatorname{Lip}(g)+C\operatorname{Lip}(\psi)$.
Comparison, followed by the reversed translation, proves
\begin{equation}
 \|v\|_{C^{0,1}(\overline\Omega)}\le C
 \label{eq:721-9}
\end{equation}
uniformly in $K,\nu$.  The Lipschitz constants of the fixed smooth extensions
of $g,\psi$ are finite.  Larger translation lengths are harmless by the
uniform value bound.  In particular the boundary normal derivative is
uniformly bounded.

\subsection{The covariance quasiderivatives and their estimates}

We use the covariance constructions in
\cite[Lemmas 6.2--6.4 and 7.4]{ZhouBellman2015}; theorem and equation numbers
here refer to the inspected arXiv version.  These constructions concern the
controlled diffusion and its defining-function barrier and have no running
payoff or discount.  The unitary adaptation requires a sign convention and
one drift term to be made explicit.

Near $\partial\Omega$, let $\eta$ be the complex column vector corresponding
to $\nabla\psi$, and let $\eta_\xi$ denote its real directional derivative.
Set
\begin{equation}
 r_0=-\frac{\operatorname{Re}(\eta^*\eta_\xi)}{|\eta|^2},
 \qquad b=-\eta_\xi-r_0\eta,
 \qquad
 Q=\frac{b\eta^*-\eta b^*}{|\eta|^2}
     -\frac{\eta^*b}{|\eta|^4}\eta\eta^*.
 \label{eq:721-10}
\end{equation}
Since $\operatorname{Re}(\eta^*b)=0$, $Q$ is skew Hermitian,
$Q\eta=b$, and $|Q|+|r_0|\le C|\xi|$.  Set
$P=-\operatorname{Real}(Q)$, the negative of its real representation.  Then
$P$ is skew symmetric, its exponential represents a unitary map, and
\begin{equation}
 D^2\psi\,\xi+r_0\nabla\psi-P\nabla\psi=0,
 \qquad
 \psi_{\xi\sigma^k}+r_0\psi_{\sigma^k}+\psi_{P\sigma^k}=0.
 \label{eq:721-11}
\end{equation}
The minus sign in $P$ is required by the column-vector convention.
This supplies the rotation allowed by the covariance cone; full orthogonal
invariance in real dimension $2n$ is not needed.  Compare the complex
construction in \cite[\S3.2, equations (3.14)--(3.15)]{ZhouHessian2013}.

There is a minor correction to the source calculation that is useful to
record.  The displayed assertion in the proof of
\cite[Lemma 7.1]{ZhouBellman2015} that
\[
 S:=\sum_k(D^2\psi\,\sigma^k,P\sigma^k)
   =\operatorname{tr}(\sigma\sigma^{\mathsf T}D^2\psi P)=0
\]
is not true in general.  Instead $|S|\le C|\xi|$.
Keeping this drift in that proof changes its term $J_4$ to
\begin{equation}
 J_4=2\frac{\psi_\xi}{\psi}
       \bigl((L\psi)_\xi+2r_0L\psi+S\bigr)
 \le C|\xi|^2+\frac{|\psi_\xi|^2}{\psi^2}.
 \label{eq:721-12}
\end{equation}
This is the same bound used there, with an enlarged geometric constant.
Taking the fixed geometric parameter $\lambda_0>0$ sufficiently small
therefore gives the same negative drift for its barrier $B_1$.
The remaining first-quasiderivative estimates use \eqref{eq:721-11}, skew symmetry,
$|P|+|r_0|\le C|\xi|$, and this negative drift; derivatives of $P$ are not
needed.  Thus the unitary version has those estimates as well.

Here is the precise input we need.  Let
$D_\delta=\{\psi>\delta\}$ and let $X_t$ be a base controlled diffusion,
starting at $x\in D_\delta$, with exit time $\tau_\delta$.
For a real vector $\xi$, the constructions supply first and second
quasiderivatives $\xi_t,\eta_t$, with $\xi_0=\xi$, $\eta_0=0$,
scalar adjoint process $Z_t$, and adjustments
$r_t,\widehat r_t,\pi_t,\widehat\pi_t,P_t$.  In the rest of the proof
$\eta_t$ denotes this second quasiderivative, independently of the gradient
vector used in \eqref{eq:721-10}.  We have
\[
 dZ_t=\pi_t\,dW_t,\qquad Z_0=0,
 \qquad \widehat\pi_t=-2\pi_tZ_t,\qquad \widehat P_t=0.
\]
The perturbed states $X_t^s$ start at $x+s\xi$ and have derivatives
$\xi_t,\eta_t$ at $s=0$.  Their density $p_t(s)$ satisfies
\begin{equation}
 p_t(0)=1,\qquad p_t'(0)=Z_t,\qquad p_t''(0)=0.
 \label{eq:721-13}
\end{equation}
The last identity follows from the choice of $\widehat\pi$ and It\^o's
formula for $Z_t^2$.  Their time factor is
\begin{equation}
 \theta_t(s)=1+\pi^{-1}\arctan\!\left(\pi(2sr_t+s^2\widehat r_t)\right),
 \quad \tfrac12<\theta_t(s)<\tfrac32,
 \quad \theta_t'(0)=2r_t,
 \quad \theta_t''(0)=2\widehat r_t.
 \label{eq:721-14}
\end{equation}
There is a positive quadratic barrier $B(x,\xi)$, bounded by $C_E|\xi|^2$
when $x$ ranges over any compact $E\Subset\Omega$, such that
\begin{equation}
 \begin{split}
 &\mathbb E\sup_{t\le\tau_\delta}(|\xi_t|^2+|Z_t|^2)
 +\mathbb E\int_0^{\tau_\delta}
 \left(|\xi_t|^2+
       \frac{|\psi_{\xi_t}(X_t)|^2}{\psi(X_t)^2}
       +r_t^2+|\widehat r_t|+|\pi_t|^2\right)dt
 \le CB(x,\xi).
 \end{split}
 \label{eq:721-15}
\end{equation}
These are \cite[Lemma 7.4(1)--(3)]{ZhouBellman2015} and the explicit
coefficient formulas.  Near the boundary
$r=r_0+\psi_\xi/\psi$ and $\widehat r=(\psi_\xi/\psi)^2$;
in the interior $\psi\ge\lambda_0^2$.
For fixed $\delta$ there is also $\ell=\ell(\delta)>0$ such that
\begin{equation}
 \sup_{\text{controls}}\sup_{\gamma\le\tau_\delta}
 \mathbb E|\xi_\gamma|^{2(1+\ell)}<\infty,
 \label{eq:721-16}
\end{equation}
by \cite[Lemma 7.4(4)]{ZhouBellman2015}.
All constants in \eqref{eq:721-15} depend only on the domain, barrier, dimension and
upper covariance bound.  In particular they are independent of $K,\nu,\delta$.

\paragraph{An additional occupation estimate and a direct second-derivative estimate.}
We shall also need
\begin{equation}
 \mathbb E\int_0^{\tau_\delta}Z_t^2dt\le CB(x,\xi).
 \label{eq:721-17}
\end{equation}
Apply It\^o's formula, with localization, to $\psi(X_t)Z_t^2$.
Using $L_A\psi\le-1$ and bounded covariances and gradient of $\psi$, its
drift is at most
\[
 -Z_t^2+C|\pi_t|^2+2C|Z_t\pi_t|
 \le-\tfrac12Z_t^2+C'|\pi_t|^2.
\]
The initial product is zero and the terminal product is nonnegative.
Expectation, followed by \eqref{eq:721-15} and monotone convergence, proves \eqref{eq:721-17}.
It follows that
\begin{equation}
 \mathbb E\int_0^{\tau_\delta}
       (|r_tZ_t|+|\pi_tZ_t|)dt\le CB(x,\xi).
 \label{eq:721-18}
\end{equation}
This is Cauchy--Schwarz on the product of probability and time measures.

The estimate for $\eta_t$ can now be proved directly, since the underlying
diffusion has zero drift.  By the source's equation (6.4),
\[
 \begin{split}
 d\eta_t&=G_t\,dW_t+H_t\,dt,\\
 G_t&=\bigl((\widehat r_t-r_t^2)I+P_t^2+2r_tP_t\bigr)\sigma_t,\\
 H_t&=-\sigma_t\widehat\pi_t-2(r_tI+P_t)\sigma_t\pi_t.
 \end{split}
\]
Put $w_t=\psi_{\xi_t}(X_t)/\psi(X_t)$.  Near the boundary the $w_t^2$
terms in $\widehat r_t-r_t^2$ cancel.  In the interior, $\psi\ge\lambda_0^2$
implies $w_t=O_{\lambda_0}(|\xi_t|)$.  Hence globally
\[
 |G_t|\le C|\xi_t|(|\xi_t|+|w_t|),\qquad
 |H_t|\le C(|\xi_t|^2+|w_t|^2)+C|\pi_tZ_t|.
\]
The Burkholder--Davis--Gundy inequality and \eqref{eq:721-15}--\eqref{eq:721-18} give
\begin{equation}
 \begin{split}
 \mathbb E\sup_{t\le\tau_\delta}|\eta_t|
 &\le C\mathbb E\left(\int_0^{\tau_\delta}|G_t|^2dt\right)^{1/2}
       +\mathbb E\int_0^{\tau_\delta}|H_t|dt\\
 &\le C\left(\mathbb E\sup|\xi_t|^2\right)^{1/2}
        \left(\mathbb E\int(|\xi_t|^2+|w_t|^2)dt\right)^{1/2}
       +CB(x,\xi)
 \le CB(x,\xi).
 \end{split}
 \label{eq:721-19}
\end{equation}
Thus no extra integrability assertion for a random exit time times a
supremum is required.

\paragraph{The unchanged terminal comparison.}
Near the boundary, the source's barrier is
\begin{equation}
 B_1(y,\zeta)=\gamma(y)
       \left(\beta(y)|\zeta|^2+
             \frac{|\psi_\zeta(y)|^2}{\psi(y)}\right),
 \qquad \gamma(y)\longrightarrow\lambda_0^2>0,
 \quad\beta(y)\longrightarrow1.
 \label{eq:721-20}
\end{equation}
Let $\gamma_s$ be the minimum of the exit times of
$X,X^s,X^{-s}$ from $D_\delta$, a bounded time localization, and
quasiderivative localizations.  For a fixed $C^2$ function $v$, the terminal
transfer in \cite[equation (9.16)]{ZhouBellman2015} gives, after taking
$s\to0$ and removing the other localizations,
\begin{equation}
 \begin{split}
 \limsup\sup_{\text{controls}}
 \mathbb E[-e^{-\phi_{\gamma_s}}
                 v_{\xi_{\gamma_s}\xi_{\gamma_s}}(X_{\gamma_s})]
 \le CB(x,\xi)
 \left(1+\sup_{\psi(y)=\delta,\,\zeta\ne0}
               \frac{|v_{\zeta\zeta}(y)|}{B_1(y,\zeta)}\right).
 \end{split}
 \label{eq:721-21}
\end{equation}
This holds for any nonnegative accumulated discount $\phi$: the terminal
comparison uses only $e^{-\phi}\le1$, and has no derivatives of $\phi$.
We spell out the needed uniform-integrability point below when applying it.

\subsection{Compensating the discount clock}

Fix $K,\nu$, so $v$ is $C^2$ up to the boundary.  Use base controls
$A_t\in\mathcal A_\nu$, $c_t\in[0,b_K]$, and put
\[
 \phi_t=\int_0^tc_sds,\qquad q_t=c_t^2/(2K)\le C_*.
\]
For the perturbed process with clock $\theta_t(s)$ choose the discount
\begin{equation}
 c_t(s)=c_t/\theta_t(s).
 \label{eq:721-22}
\end{equation}
This is admissible for the larger cap because $\theta>1/2$.
The covariance rotations are admissible by unitary invariance.  The identities
\begin{equation}
 \theta_t(s)c_t(s)=c_t,
 \qquad
 \theta_t(s)\left(c_t(s)g(X_t^s)-\frac{c_t(s)^2}{2K}\right)
 =c_tg(X_t^s)-\frac{q_t}{\theta_t(s)}
 \label{eq:721-23}
\end{equation}
show that the killing clock is exactly independent of $s$.

For the common localized stopping time $\gamma_s$, set
\begin{equation}
 J_s=\mathbb E\left[
 e^{-\phi_{\gamma_s}}p_{\gamma_s}(s)v(X_{\gamma_s}^s)
 +\int_0^{\gamma_s}e^{-\phi_t}p_t(s)
       \left(c_tg(X_t^s)-\frac{q_t}{\theta_t(s)}\right)dt
 \right].
 \label{eq:721-24}
\end{equation}
Then $J_s\le v(x+s\xi)$.  Here is a direct verification, avoiding any
filtration issue in the transformed control.  In It\^o's formula for
$e^{-\phi_t}p_t(s)v(X_t^s)$, the Girsanov drift in $X^s$ cancels the
quadratic-covariation term with $p(s)$.  After adding the integral in
\eqref{eq:721-24}, the remaining drift divided by $e^{-\phi_t}p_t(s)>0$ is
\[
 \theta_t(s)\left[
 L_{A_t(s)}v-c_t(s)v+c_t(s)g-\frac{c_t(s)^2}{2K}
 \right]\le0,
\]
where $A_t(s)$ is the admissible unitary conjugate of $A_t$.
This is precisely the larger-cap Bellman inequality.
The bounded localizations make the stochastic integrals true martingales.

At $s=0$, the supremum of $J_0$ over smaller-cap base controls is $v(x)$,
by the usual Bellman verification/dynamic-programming equality and the
coincidence of the two value functions.  This remains true for the common
control-dependent bounded stopping time; an approximately optimal policy
for the smaller-cap equation may be stopped at that time.
Consequently
\begin{equation}
 -\frac{v(x+s\xi)-2v(x)+v(x-s\xi)}{s^2}
 \le\sup_{\text{base controls}}
       \frac{2J_0-J_s-J_{-s}}{s^2}.
 \label{eq:721-25}
\end{equation}
For each base control the same $\gamma_s$ is used in all three terms.  There
is therefore no differentiation of an indicator of an obstacle-stopping
region.  The value function being differentiated is the smooth approximant.

\subsection{Every changed second-variation term is uniformly bounded}

First take $s\to0$ with the time and quasiderivative localizations fixed.
The finite-time convergence in \cite[Lemmas 6.2--6.3]{ZhouBellman2015}
justifies the state and density expansions uniformly over the base controls.
For each fixed $K$ the cap $b_K$ is finite.  Remainder convergence may depend
on $K,\nu$, while the estimate obtained after taking this limit will not.

By \eqref{eq:721-13} and the constant killing clock, the terminal second variation is
\begin{equation}
 e^{-\phi_\gamma}\left[
 v_{\xi_\gamma\xi_\gamma}(X_\gamma)
 +v_{\eta_\gamma}(X_\gamma)
 +2Z_\gamma v_{\xi_\gamma}(X_\gamma)
 \right].
 \label{eq:721-26}
\end{equation}
The last two terms have expected absolute value at most $CB(x,\xi)$ by
\eqref{eq:721-9}, \eqref{eq:721-15}, \eqref{eq:721-19}, and Cauchy--Schwarz.

The running second variation is exactly
\begin{equation}
 \begin{split}
 R_\gamma=\int_0^\gamma e^{-\phi_t}\bigl[
 &c_t\{g_{\xi_t\xi_t}(X_t)+g_{\eta_t}(X_t)
                   +2Z_tg_{\xi_t}(X_t)\}\\
 &-q_t\{8r_t^2-2\widehat r_t-4r_tZ_t\}\bigr]dt.
 \end{split}
 \label{eq:721-27}
\end{equation}
Indeed $(1/\theta)'(0)=-2r$ and
$(1/\theta)''(0)=8r^2-2\widehat r$, while $p'(0)=Z$ and $p''(0)=0$.
This accounts for all terms introduced by the quadratic penalty.

For the terms with $c_t$, use the pathwise identity
\[
 \int_0^\gamma c_te^{-\phi_t}dt\le1.
\]
Their expected absolute value is at most
\[
 C\|g\|_{C^2}\mathbb E\left[
 \sup|\xi_t|^2+\sup|\eta_t|+2\sup|Z_t\xi_t|
 \right]\le CB(x,\xi).
\]
For the other terms, use $q_t\le C_*$, \eqref{eq:721-15}, and \eqref{eq:721-18}.
Thus
\begin{equation}
 \mathbb E|R_\gamma|\le CB(x,\xi),
 \label{eq:721-28}
\end{equation}
uniformly in $K,\nu$ and every localization.

For clarity, the passage through the localizations uses the following precise
facts.  Keep $\delta,K,\nu$ fixed.  Stop at time $T$ and when
$|\xi_t|+|Z_t|$ first exceeds $N$.  Then $r,\widehat r,P,\pi,\widehat\pi$
are bounded, and the finite-time difference-quotient lemmas apply to the
common stopping time.  After $s\to0$, \cite[Lemma 6.4]{ZhouBellman2015}
controls the difference between the base exit time and the minimum of the
three exit times.  Formula \eqref{eq:721-16} makes $|\xi_\gamma|^2$ uniformly
integrable over controls and stopped times.  The probability of the
$N$-cutoff is at most $CB/N^2$ by \eqref{eq:721-15}, while
$\mathbb P(\tau_\delta>T)\le\|\psi\|_\infty/T$ by \eqref{eq:721-3}.
H\"older's inequality therefore removes these exceptional events from the
terminal Hessian term, bounded for the fixed approximant by
$\|D^2v\|_\infty|\xi_\gamma|^2$.
On the remaining events uniform continuity of $D^2v$ on the fixed closed
set $\overline D_\delta$, together with the exit-time comparison, gives
\eqref{eq:721-21}.  The intermediate continuity moduli and higher-moment constants
may depend on $\delta,K,\nu$; their sole role is to make errors tend to zero
before these parameters vary.  The lower-order terminal and running terms
need no convergence statement: their absolute expectations are bounded by
$CB$ at every localization, so their bounds survive a limsup.

Combining \eqref{eq:721-25}--\eqref{eq:721-28} in this order yields
\begin{equation}
 -v_{\xi\xi}(x)\le CB(x,\xi)
 +\limsup\sup_{\text{controls}}
       \mathbb E[-e^{-\phi_\gamma}
                         v_{\xi_\gamma\xi_\gamma}(X_\gamma)].
 \label{eq:721-29}
\end{equation}
The remaining term is the unchanged terminal term in \eqref{eq:721-21}.

\subsection{Boundary passage and all real second derivatives}

For each fixed $K,\nu$, $v\in C^2(\overline\Omega)$.
If $y\in\partial\Omega$ and $\zeta$ is real tangent to the boundary,
differentiating $v=g$ along a boundary curve gives
\begin{equation}
 |v_{\zeta\zeta}(y)|
 \le C\left(\|g\|_{C^2}
                  +\|\nabla v\|_{L^\infty(\partial\Omega)}\right)
          |\zeta|^2\le C|\zeta|^2.
 \label{eq:721-30}
\end{equation}
This includes the second fundamental form term; Euclidean tangential second
derivatives are not identified with those of an arbitrary extension of $g$.

By \eqref{eq:721-20}, every limiting nontangential direction has a denominator tending
to infinity.  Limiting tangent directions have numerator bounded by
\eqref{eq:721-30}.  For fixed $K,\nu$, continuity of $D^2v$ makes the normal and
mixed contributions vanish as the directions approach a tangent direction.
Normalize $|\zeta|=1$ and use compactness to obtain
\begin{equation}
 \limsup_{\delta\downarrow0}
 \sup_{\psi(y)=\delta,\,\zeta\ne0}
       \frac{|v_{\zeta\zeta}(y)|}{B_1(y,\zeta)}\le C,
 \label{eq:721-31}
\end{equation}
where $C$ is independent of $K,\nu$.
The order matters: reach the boundary for each fixed smooth approximant
before varying the approximation parameters.  Its possibly large normal
Hessian disappears from this ratio.

Equations \eqref{eq:721-21}, \eqref{eq:721-29}, and \eqref{eq:721-31} give
\begin{equation}
 v_{\xi\xi}(x)\ge-CB(x,\xi).
 \label{eq:721-32}
\end{equation}
On each compact $E\Subset\Omega$, this is a uniform lower bound
$D^2v\ge-C_EI$.  Moreover \eqref{eq:721-6} implies $F_\nu(D^2v)\le0$.
Since $I/n\in\mathcal A_\nu$, $\Delta_{\mathbb R^{2n}}v\le0$.
The lower eigenvalue bound and this nonpositive trace bound each real
Hessian eigenvalue above by $(2n-1)C_E$.  We have proved
\begin{equation}
 \|v_{K,\nu}\|_{C^{1,1}(E)}\le C_E
 \label{eq:721-33}
\end{equation}
with constants depending on $\Omega,\rho,m,n,E$ and a $C^2$ bound for $h$,
but not on $K,\nu$.

\subsection{Identification of the limit with the exact envelope}

Choose smooth $h_j\to h$ in $C^2(\overline\Omega)$,
$K_j\to\infty$, $\nu_j\to0$, and put $u_j=-v_{K_j,\nu_j}$ with
$g_j=-h_j$.  The preceding uniform estimates give a subsequence converging
uniformly on $\overline\Omega$ and in $C^1$ on compact subsets to
\[
 u\in C^{0,1}(\overline\Omega)\cap C^{1,1}_{\mathrm{loc}}(\Omega).
\]
By \eqref{eq:721-6}, $u_j\le h_j+C/\sqrt{K_j}$, so $u\le h$.
For each $A\in\mathcal A$ set
$A_j=(1-\nu_j)A+\nu_jI/n$.  Formula \eqref{eq:721-6} gives $L_{A_j}u_j\ge0$.
Passing to distributions shows $L_Au\ge0$ for all $A$, hence $u$ is
$m$-subharmonic.  The boundary barriers give $u=h$ on $\partial\Omega$.

Let $w\in\operatorname{SH}_m(\Omega)$ with $w\le h$ be any competitor.
A fixed sufficiently large $C$ makes $h+C\rho$ $m$-subharmonic, by the
uniform strictness of $\rho$.  Replacing $w$ by
$\max(w,h+C\rho)$ increases it and allows us to assume it is bounded.
Put $\varepsilon_j=\|h_j-h\|_\infty$.  For $t>0$ the function
\[
 w_{j,t}=w-\varepsilon_j+t\rho
\]
is strictly subharmonic for every $L_{A_j}$, is at most $h_j$, and has
boundary upper limits at most $h_j$.  If $w_{j,t}-u_j$ had a positive
interior maximum, the smooth function $u_j$ plus a constant would touch
$w_{j,t}$ from above there.  The linear viscosity inequalities would imply
$L_{A_j}u_j\ge t\delta_0$ for every $A_j$, hence
$F_{\nu_j}(D^2v_j)\le-t\delta_0$.  On the other hand, at this point
$u_j<w_{j,t}\le h_j$, so the penalty vanishes and \eqref{eq:721-6} gives
$F_{\nu_j}(D^2v_j)=0$, a contradiction.  Thus $w_{j,t}\le u_j$.
Let $t\downarrow0$ and $j\to\infty$ to obtain $w\le u$.

The function $u$ is an admissible competitor and dominates every admissible
competitor.  Therefore $u=P_{m,\Omega}h$, and \eqref{eq:721-33} proves the theorem.

\begingroup\small
\begin{thebibliography}{99}
\bibitem{BenaliZeriahi} A. Benali and A. Zeriahi, \emph{The H\"older continuous subsolution theorem for complex Hessian equations}, J. \'Ec. polytech. Math. \textbf{7} (2020), 981--1007, p.~998. \url{https://doi.org/10.5802/jep.133}.
\bibitem{BenaliErratum} A. Benali and A. Zeriahi, \emph{Erratum to The H\"older continuous subsolution theorem for complex Hessian equations}, J. \'Ec. polytech. Math. \textbf{8} (2021), 779--789. \url{https://doi.org/10.5802/jep.157}.
\bibitem{ZhouBellman2015} W. Zhou, \emph{Interior regularity of fully nonlinear degenerate elliptic equations I: Bellman equations with constant coefficients}, SIAM J. Math. Anal. \textbf{47} (2015), 2375--2415. \url{https://doi.org/10.1137/130912724}. Numbered passages cited in the text refer to the inspected preprint \url{https://arxiv.org/abs/1302.7062v2}.
\bibitem{Krylov1987} N. V. Krylov, \emph{Nonlinear Elliptic and Parabolic Equations of the Second Order}, Mathematics and its Applications 7, D. Reidel, Dordrecht, 1987, \S6.2, Theorem 7, as quoted in Zhou's Remark 5.1. \url{https://link.springer.com/book/9781402003349}.
\bibitem{ZhouHessian2013} W. Zhou, \emph{Representation and regularity for the Dirichlet problem for real and complex degenerate Hessian equations}, preprint (2013), Theorem 6.1. \url{https://arxiv.org/abs/1311.6450}.
\end{thebibliography}
\endgroup
\end{document}
